% ---------------------------------------------------------------
% tables.tex -- generated by tables_fixed.py from the CORRECTED run
%   label: 50% KDE core-area volume contour (was: median of nonzero
%          on a linear-interpolated surface)
%   CV   : contiguous leave-one-region-out, k-means (was: GroupKFold
%          over 3-degree blocks, which scattered them)
% Requires: booktabs, multirow
% ---------------------------------------------------------------


\begin{table}[t]
\centering
\caption{TabICLv2 versus a tuned Random Forest under \textbf{contiguous leave-one-region-out} cross-validation, with the \textbf{KDE core-area} label. Both models are scored on identical class-balanced test cells per fold. Differences are paired per fold; intervals are 95\% $t$-intervals on $n-1$ degrees of freedom.}
\label{tab:parity-caretta}
\begin{tabular}{lccccc}
\toprule
Metric & TabICLv2 & RF-full & Paired $\Delta$ & 95\% CI & $p$ \\
\midrule
AUC & 0.9878 & 0.9089 & +0.0789 & $[-0.088, +0.246]$ & 0.231 \\
F1 & 0.8689 & 0.7924 & +0.0765 & $[+0.027, +0.126]$ & \textbf{0.016} \\
\bottomrule
\end{tabular}
\end{table}

\begin{table}[t]
\centering
\caption{Robustness across the held-out regions. Variance ratios are RF/TabICLv2 on the squared standard deviations; range and worst-fold rows compare the values shown directly.}
\label{tab:robustness-caretta}
\begin{tabular}{lccc}
\toprule
Robustness summary & TabICLv2 & RF-full & Ratio \\
\midrule
AUC sd across folds & 0.018 & 0.123 & 45.8$\times$ lower variance \\
F1 sd across folds & 0.137 & 0.165 & 1.4$\times$ lower variance \\
AUC fold range & 0.038 & 0.256 & 6.7$\times$ narrower \\
F1 fold range & 0.304 & 0.360 & 1.2$\times$ narrower \\
Worst-fold AUC & 0.961 & 0.724 & +0.236 \\
Worst-fold F1 & 0.673 & 0.551 & +0.122 \\
\bottomrule
\end{tabular}
\end{table}

\begin{table}[t]
\centering
\caption{Model performance by species (mean, standard deviation, and worst-case AUC and F1) under the KDE core-area label and contiguous leave-one-region-out CV.}
\label{tab:by-species}
\begin{tabular}{llcccccc}
\toprule
\multirow{2}{*}{Species} & \multirow{2}{*}{Model} & \multicolumn{3}{c}{AUC} & \multicolumn{3}{c}{F1} \\
\cmidrule(lr){3-5}\cmidrule(lr){6-8}
 & & Mean & SD & Worst & Mean & SD & Worst \\
\midrule
\emph{Caretta caretta} & TabICLv2 & 0.9878 & 0.0182 & 0.9608 & 0.8689 & 0.1369 & 0.6726 \\
 & RF-full & 0.9089 & 0.1233 & 0.7244 & 0.7924 & 0.1647 & 0.5507 \\
 & AquaMaps$^{\dagger}$ & 0.5728 & 0.1274 & 0.4308 & 0.2708 & 0.2931 & 0.0247 \\
\addlinespace
\emph{Cetorhinus maximus} & TabICLv2 & 0.8543 & 0.3018 & 0.3149 & 0.7861 & 0.3673 & 0.1366 \\
 & RF-full & 0.8459 & 0.2742 & 0.3588 & 0.7469 & 0.3662 & 0.1006 \\
 & AquaMaps$^{\dagger}$ & 0.5513 & 0.0744 & 0.5000 & 0.6667 & 0.0000 & 0.6667 \\
\addlinespace
\emph{Balaenoptera musculus} & TabICLv2 & 0.9203 & 0.0596 & 0.8848 & 0.7022 & 0.1669 & 0.5094 \\
 & RF-full & 0.8964 & 0.0582 & 0.8544 & 0.4710 & 0.4196 & 0.1227 \\
 & AquaMaps$^{\dagger}$ & 0.4609 & 0.1792 & 0.2621 & 0.6287 & 0.0309 & 0.6045 \\
\addlinespace
\bottomrule
\end{tabular}
\begin{minipage}{\linewidth}\footnotesize $^{\dagger}$AquaMaps is the published species-distribution surface. It is \emph{not} fitted to this data -- a fixed 0.5$^{\circ}$ envelope re-scored on the identical test cells -- so its across-fold spread reflects only which region was held out, and no paired test against it is reported. The comparison is partly circular in its favour: its occurrence base overlaps OBIS-SEAMAP.\end{minipage}
\end{table}

\begin{table}[t]
\centering
\caption{TabICLv2 against RF-full under contiguous leave-one-region-out cross-validation. $\Delta$ is TabICLv2 minus RF-full, so a positive value favours TabICLv2.}
\label{tab:parity-all}
\begin{tabular}{lccccc}
\toprule
Metric & TabICLv2 & RF-full & $\Delta$ & $p$ & Folds won \\
\midrule
\multicolumn{6}{l}{\emph{Caretta caretta}, Mediterranean} \\
\quad AUC & 0.9878 & 0.9089 & +0.0789 & 0.231 & 4/4 \\
\quad F1 & 0.8689 & 0.7924 & +0.0765 & \textbf{0.016} & 4/4 \\
\addlinespace
\multicolumn{6}{l}{\emph{Cetorhinus maximus}, NE Atlantic} \\
\quad AUC & 0.8543 & 0.8459 & +0.0085 & 0.623 & 3/5 \\
\quad F1 & 0.7861 & 0.7469 & +0.0392 & 0.096 & 4/5 \\
\addlinespace
\multicolumn{6}{l}{\emph{Balaenoptera musculus}, California Current} \\
\quad AUC & 0.9203 & 0.8964 & +0.0238 & \textbf{0.036} & 3/3 \\
\quad F1 & 0.7022 & 0.4710 & +0.2311 & 0.434 & 2/3 \\
\addlinespace
\bottomrule
\end{tabular}
\begin{minipage}{\linewidth}\footnotesize Bold $p$ values are significant at $\alpha=0.05$. Folds won counts the held-out regions on which TabICLv2 scored higher than RF-full.\end{minipage}
\end{table}

\begin{table}[t]
\centering
\caption{Dispersion of scores across the held-out regions, and the weakest region for each model.}
\label{tab:dispersion}
\begin{tabular}{lccccccc}
\toprule
 & \multicolumn{2}{c}{Across-fold sd} & \multicolumn{3}{c}{Variance ratio, RF/TabICLv2} & \multicolumn{2}{c}{Worst fold} \\
\cmidrule(lr){2-3}\cmidrule(lr){4-6}\cmidrule(lr){7-8}
Metric & TabICLv2 & RF-full & Ratio & $p_F$ & $p_{\text{Levene}}$ & TabICLv2 & RF-full \\
\midrule
\multicolumn{8}{l}{\emph{Caretta caretta}, Mediterranean} \\
\quad AUC & 0.018 & 0.123 & 45.81 & \textbf{0.011} & 0.382 & 0.961 & 0.724 \\
\quad F1 & 0.137 & 0.165 & 1.45 & 0.769 & 0.903 & 0.673 & 0.551 \\
\addlinespace
\multicolumn{8}{l}{\emph{Cetorhinus maximus}, NE Atlantic} \\
\quad AUC & 0.302 & 0.274 & 0.83 & 1.000 & 0.999 & 0.315 & 0.359 \\
\quad F1 & 0.367 & 0.366 & 0.99 & 1.000 & 0.984 & 0.137 & 0.101 \\
\addlinespace
\multicolumn{8}{l}{\emph{Balaenoptera musculus}, California Current} \\
\quad AUC & 0.060 & 0.058 & 0.95 & 1.000 & 0.976 & 0.885 & 0.854 \\
\quad F1 & 0.167 & 0.420 & 6.32 & 0.273 & 0.420 & 0.509 & 0.123 \\
\addlinespace
\bottomrule
\end{tabular}
\begin{minipage}{\linewidth}\footnotesize Variance ratios are RF-full over TabICLv2, so a ratio below 1 means TabICLv2 was the more variable of the two. Bold entries are significant at $\alpha=0.05$.\end{minipage}
\end{table}

\begin{table}[t]
\centering
\caption{Size and class balance of each task under the corrected pipeline. Prevalence is the fraction of ocean cells inside the 50\% KDE core-area contour.}
\label{tab:task-scale}
\begin{tabular}{lccc}
\toprule
 & \emph{C. caretta} & \emph{C. maximus} & \emph{B. musculus} \\
 & Mediterranean & NE Atlantic & California Current \\
\midrule
Grid cells & 13,683 & 36,902 & 42,476 \\
Prevalence (core area) & 0.292 & 0.219 & 0.052 \\
Prevalence (old label) & 0.435 & 0.263 & 0.129 \\
Context rows per fold & 10,891 & 29,521 & 33,785 \\
Test cells per fold & 300 & 253 & 300 \\
Regions scored / attempted & 4 / 5 & 5 / 5 & 3 / 5 \\
Train--test separation (median, $^\circ$) & 2.40 & 1.93 & 3.80 \\
Occurrence rows after filtering & 228 & 142 & 41 \\
\bottomrule
\end{tabular}
\end{table}

\begin{table}[t]
\centering
\caption{Language-model form of in-context learning, and sensitivity checks, under the corrected pipeline. TabLLM is Qwen2.5-72B-Instruct-AWQ scored by label-verbalizer logprob.}
\label{tab:tabllm}
\begin{tabular}{ll}
\toprule
Quantity & Value \\
\midrule
TabICLv2 against TabLLM, matched $k=32$ & 2/3 species; $\Delta$AUC +0.306, +0.129, -0.001 \\
TabLLM 72B AUC (mean per species) & 0.555, 0.660, 0.893 \\
TabLLM 72B AUC (best / worst fold-draw) & 0.989 / 0.375 \\
Truncation-fabricated $P(\text{HIGH})=0$ & 0 of 16,830 predictions \\
Exemplar-draw variance ($k$-shot) & up to 0.335 AUC swing within a fold \\
\bottomrule
\end{tabular}
\begin{minipage}{\linewidth}\footnotesize $\Delta$AUC is TabICLv2 minus TabLLM at matched $k$, in species order. Species are listed in the order of Table~\ref{tab:by-species}. Logprob scoring emits a single verbalizer token, so the truncation failure that fabricated confident zeros in the earlier verbalized run cannot occur.\end{minipage}
\end{table}

\begin{table}[t]
\centering
\caption{Permutation importance on the contiguous spatial holdout, as the mean drop in AUC when each feature is permuted (RF-full, averaged over held-out regions).}
\label{tab:perm-importance}
\begin{tabular}{lccc}
\toprule
Feature & \emph{C. caretta} & \emph{C. maximus} & \emph{B. musculus} \\
\midrule
Latitude & 0.367 & 0.036 & -0.000 \\
Longitude & 0.015 & 0.093 & 0.038 \\
Distance to coast & 0.008 & 0.077 & 0.282 \\
Speed & 0.008 & 0.010 & 0.003 \\
Retention & 0.002 & 0.012 & -0.002 \\
Vorticity & 0.001 & 0.003 & -0.007 \\
\bottomrule
\end{tabular}
\end{table}

\begin{table}[t]
\centering
\caption{Zero-shot against fine-tuned TabICLv2. \textbf{Not re-run under the corrected pipeline.}}
\label{tab:finetune}
\begin{tabular}{lcccc}
\toprule
Species & ZS AUC & FT AUC & $\Delta$ & Folds won \\
\midrule
\emph{Caretta caretta} & \emph{not run} & \emph{not run} & \emph{not run} & \emph{not run} \\
\emph{Cetorhinus maximus} & \emph{not run} & \emph{not run} & \emph{not run} & \emph{not run} \\
\emph{Balaenoptera musculus} & \emph{not run} & \emph{not run} & \emph{not run} & \emph{not run} \\
\bottomrule
\end{tabular}
\begin{minipage}{\linewidth}\footnotesize The fine-tuning study runs on Modal and bills GPU credits; it was not relaunched. Its published values were produced under the superseded label and CV scheme and are not carried forward.\end{minipage}
\end{table}

\begin{table}[t]
\centering
\caption{ROC-AUC for full-context and $k=32$ models, averaged over held-out regions. $k$-shot rows are additionally averaged over exemplar draws; $\pm$ is one standard deviation across regions.}
\label{tab:models-k32}
\begin{tabular}{lccc}
\toprule
Model & \emph{Caretta caretta} & \emph{Cetorhinus maximus} & \emph{Balaenoptera musculus} \\
\midrule
TabICLv2 (full context) & 0.988 $\pm$ 0.018 & 0.854 $\pm$ 0.302 & 0.920 $\pm$ 0.060 \\
RF-full (full context) & 0.909 $\pm$ 0.123 & 0.846 $\pm$ 0.274 & 0.896 $\pm$ 0.058 \\
\addlinespace
TabICLv2 ($k=32$) & 0.862 $\pm$ 0.071 & 0.789 $\pm$ 0.245 & 0.892 $\pm$ 0.045 \\
TabLLM (Qwen2.5-72B) ($k=32$) & 0.555 $\pm$ 0.125 & 0.660 $\pm$ 0.133 & 0.893 $\pm$ 0.064 \\
Random Forest ($k=32$) & 0.822 $\pm$ 0.095 & 0.802 $\pm$ 0.216 & 0.859 $\pm$ 0.073 \\
GBT ($k=32$) & 0.835 $\pm$ 0.126 & 0.721 $\pm$ 0.212 & 0.789 $\pm$ 0.118 \\
kNN ($k=32$) & 0.656 $\pm$ 0.099 & 0.700 $\pm$ 0.112 & 0.729 $\pm$ 0.082 \\
LogReg ($k=32$) & 0.374 $\pm$ 0.110 & 0.423 $\pm$ 0.159 & 0.793 $\pm$ 0.099 \\
\bottomrule
\end{tabular}
\begin{minipage}{\linewidth}\footnotesize All $k=32$ rows draw the \emph{same} 32 exemplars and are scored on the \emph{same} test cells. TabLLM is Qwen2.5-72B-Instruct-AWQ under logprob scoring (the label-verbalizer likelihood the method specifies), served on 4$\times$A10G.\end{minipage}
\end{table}
